\chapter{Fidélité généalogique du pont matriciel : domaine, correspondance et limites exactes}
\label{gmc:app:certificado}

\section{Identités de naturalité, de positivité et de compatibilité du pont généalogique–matriciel}

Les contrôles suivants s'obtiennent en reconstruisant les opérateurs à partir de leurs
données d'incidence. Les identités rationnelles sont calculées exactement ; les
comparaisons en virgule flottante déclarent leur borne d'erreur.

\begin{longtable}{@{}L{.35\textwidth}L{.22\textwidth}L{.32\textwidth}@{}}
\toprule
Contrôle & Statut & Erreur maximale ou identité\\
\midrule
\endfirsthead
\toprule
Contrôle & Statut & Erreur maximale ou identité\\
\midrule
\endhead
cuatro \PVM qutrit & vérifié & \(4.72\times10^{-16}\)\\
idempotence des douze projecteurs & vérifiée & \(3.34\times10^{-16}\)\\
solapamientos MUB & vérifiés & \(\lvert\Tr(PQ)-1/3\rvert<3.34\times10^{-16}\)\\
obstruction de colonne QLS & exacte & minimum \(1/3\)\\
QMS aditivo \((9,3)\) & vérifié & lignes et colonnes exactes à la tolérance près\\
QMS à trois contextes \((9,3)\) & vérifié & erreur \(<4.21\times10^{-16}\)\\
non-commutativité des cellules & vérifiée & norme maximale \(0.074074\ldots\)\\
QMS atlas \((12,3)\) & vérifié & erreur \(<4.24\times10^{-16}\)\\
décomposition semi-classique & exacte & erreur de reconstruction nulle\\
cube qutrit d'ordre neuf & vérifié & erreur de ligne \(<4.21\times10^{-16}\)\\
cube qutrit d'ordre douze & vérifié & erreur de ligne \(<4.24\times10^{-16}\)\\
\(\mathbb P^1(\mathbb Z/9)\to\mathbb P^1(\mathbb F_3)\) & exact &
\(12\to4\) con fibras \((3,3,3,3)\)\\
troncature--compression & exacte & conmutador cero\\
lecteur direct \(1/\pi\) & exact & normalisation unitaire\\
calibre MUB \(3/\pi\) & vérifié & probabilidad \(1/\pi\), error \(<3.85\times10^{-16}\)\\
polynôme de Gram APP & exact & \(t^5(t-1)(t-5/7)(t-1/3)^2\)\\
quatre intersections de Halmos & exactes & \((64,1,5,5)\)\\
algèbre \(C^*\) d'APP dans \(d=2\) & exacte &
\(\mathbb C^4\oplus M_2\oplus M_2\), dimension \(12\)\\
\bottomrule
\end{longtable}

La tolérance des contrôles numériques est \(2\times10^{-10}\). Les
énoncés rationnels n'utilisent pas de tolérance : ils sont obtenus par fractions exactes
et élimination de Gauss sur \(\mathbb Q\).

\Needspace{30\baselineskip}
\section{Coefficients exacts du polynôme}

La construction produit
\[
 \det(tI-CC^*)
 =t^9-\frac{50}{21}t^8+\frac{124}{63}t^7
  -\frac23t^6+\frac5{63}t^5,
 \tag{GMC-A.1}
\]
et coïncide coefficient par coefficient avec
\[
 t^5(t-1)(t-5/7)(t-1/3)^2.
 \tag{GMC-A.2}
\]
La vérification construit \(CC^*\) à partir des cardinalités exactes des
intersections de fibres et applique Faddeev--LeVerrier sur \(\mathbb Q\) ; elle ne
dépend pas d'une diagonalisation numérique.

\section{Obstructions à l'identification d'équivalences non démontrées dans le pont matriciel}

Le certificat fini ne remplace ni la filtration nonadique à profondeur arbitraire, ni la
Loi générale des masses, ni les dérivations de Planck. Il ne démontre pas non plus à lui
seul la fidélité APP--TRIT--TPK de tout QMS cyclique, l'équivariance complète
de la calibration \(\Xi\), l'identification de la mémoire TPK à un système
auxiliaire de Stinespring, la divisibilité infinie d'un canal MPS ni une
universalité catégorique. Chacun de ces énoncés conserve son emplacement
principal et ses propres réfutateurs.
